\chapter{Putting It All Together}

A small editor-style application is the best way to practice \api{crystal-qt6}. It uses almost every desktop concept from the earlier chapters: a main window, actions, menus, toolbars, docks, model/view panels, dialogs, custom painting, pointer events, widget grabs, and status messages.

The maintained reference for this chapter is \api{examples/editor_vertical_slice.cr}. It is intentionally smaller than a full editor, but it is complete enough to run, pan, zoom, switch layers, adjust inspector values, choose a color, and export a PNG snapshot.

\section{What The Slice Builds}

The vertical slice has four visible regions:

\begin{itemize}
\item a central canvas implemented with \api{EventWidget};
\item a left \api{DockWidget} containing a model-backed layer tree;
\item a right \api{DockWidget} containing inspector controls;
\item a menu bar, toolbar, and status bar for application commands.
\end{itemize}

The goal is not to build every feature at once. The goal is to make one thin path work from end to end: select a layer, redraw the canvas, adjust state, and export the result. Once that loop is reliable, larger features have somewhere sensible to attach.

\begin{figure}[tbp]
  \centering
  \screenshotimage{docs/book/images/putting-together-editor-shell.png}{width=0.95\textwidth,height=0.42\textheight,keepaspectratio}%
  \caption{The vertical slice combines a central canvas, model/view layer dock, inspector dock, toolbar, menus, and status bar.}
\end{figure}

\begin{figure}[tbp]
  \centering
  \screenshotimage{docs/book/images/putting-together-architecture-map.png}{width=0.95\textwidth,height=0.46\textheight,keepaspectratio}%
  \caption{Source map for \api{examples/editor_vertical_slice.cr}; numbered regions connect the running UI to the relevant code.}
\end{figure}

\section{Start With Application State}

Keep document-ish state separate from the widgets. In the example, \api{EditorVerticalSliceState} stores the current layer, zoom, pan, grid spacing, marker size, accent color, and drag state.

\begin{lstlisting}[style=crystal]
class EditorVerticalSliceState
  property active_layer : String
  property zoom : Float64
  property pan_x : Float64
  property pan_y : Float64
  property grid_spacing : Int32
  property marker_size : Int32
  property show_grid : Bool
  property accent : Qt6::Color

  def initialize
    @active_layer = "Terrain"
    @zoom = 1.0
    @pan_x = 32.0
    @pan_y = 36.0
    @grid_spacing = 48
    @marker_size = 18
    @show_grid = true
    @accent = Qt6::Color.new(62, 130, 109)
  end
end
\end{lstlisting}

This object is deliberately ordinary Crystal code. It can be tested, serialized, or replaced later without rewriting the UI. Widgets read from it and write to it, but they do not become the document model.

\section{Create The Shell}

Start with a \api{MainWindow}, a status bar, and a central canvas. Leave the docks and actions until the central workflow can draw something.

\begin{lstlisting}[style=crystal]
app = Qt6.application
state = EditorVerticalSliceState.new

main = Qt6::MainWindow.new
main.window_title = "Editor Vertical Slice"
main.resize(1240, 780)

status_bar = main.status_bar
status_bar.show_message("Ready")

canvas = Qt6::EventWidget.new
canvas.resize(860, 620)
main.central_widget = canvas
\end{lstlisting}

This shell is small, but it already establishes ownership. The main window owns the central widget, status bar, docks, menus, and toolbars that will be added later.

\section{Draw The Canvas}

The canvas uses \api{EventWidget#on_paint_with_painter}. The callback receives a live paint event and a live \api{QPainter}; use them only inside the callback.

\begin{lstlisting}[style=crystal]
scene_to_view = ->(point : Qt6::PointF) do
  Qt6::PointF.new(
    state.pan_x + point.x * state.zoom,
    state.pan_y + point.y * state.zoom
  )
end

canvas.on_paint_with_painter do |event, painter|
  painter.fill_rect(event.rect, Qt6::Color.new(245, 243, 238))

  scene_rect = Qt6::RectF.new(0.0, 0.0, 640.0, 420.0)
  visible_rect = Qt6::RectF.new(
    state.pan_x,
    state.pan_y,
    scene_rect.width * state.zoom,
    scene_rect.height * state.zoom
  )

  painter.pen = Qt6::QPen.new(Qt6::Color.new(74, 82, 91), 2.0)
  painter.brush = Qt6::Color.new(255, 255, 255)
  painter.draw_rect(visible_rect)
end
\end{lstlisting}

The real example adds a grid, a route, markers, and a small heads-up display. The important shape is the same: read state, transform scene coordinates into view coordinates, draw, and return. When state changes, call \api{canvas.update}; do not paint from mouse or wheel callbacks.

\begin{figure}[tbp]
  \centering
  \screenshotimage{docs/book/images/putting-together-canvas.png}{width=0.84\textwidth,height=0.44\textheight,keepaspectratio}%
  \caption{The central \api{EventWidget} paints the current layer from ordinary application state.}
\end{figure}

\section{Add Pointer Interaction}

Pointer callbacks mutate state and schedule a repaint. The pan interaction needs only the last pointer position and a dragging flag.

\begin{lstlisting}[style=crystal]
canvas.on_mouse_press do |event|
  state.dragging = true
  state.last_pointer = event.position
  status_bar.show_message("Dragging canvas", 1000)
end

canvas.on_mouse_move do |event|
  next unless state.dragging

  state.pan_x += event.position.x - state.last_pointer.x
  state.pan_y += event.position.y - state.last_pointer.y
  state.last_pointer = event.position
  canvas.update
end

canvas.on_mouse_release do |_event|
  state.dragging = false
  status_bar.show_message("Canvas settled", 900)
end
\end{lstlisting}

Wheel and key callbacks follow the same rule. Change the zoom, clamp it to a reasonable range, update the canvas, and show a short status message.

\section{Add A Layer Model}

The layer panel is model/view rather than a hand-managed list. The source model stores names, states, priorities, and colors. A proxy sorts the layers by priority.

\begin{lstlisting}[style=crystal]
layer_model = Qt6::StandardItemModel.new(main)
layer_model.set_horizontal_header_label(0, "Layer")
layer_model.set_horizontal_header_label(1, "State")

[
  {"Terrain", "Visible", 10, Qt6::Color.new(62, 130, 109)},
  {"Units", "Visible", 20, Qt6::Color.new(204, 86, 62)},
  {"Labels", "Draft", 30, Qt6::Color.new(92, 88, 176)},
].each_with_index do |entry, row|
  name = Qt6::StandardItem.new(entry[0])
  name.set_data(entry[2], Qt6::ItemDataRole::User)
  name.set_data(entry[3], Qt6::ItemDataRole::Foreground)

  layer_model.set_item(row, 0, name)
  layer_model.set_item(row, 1, Qt6::StandardItem.new(entry[1]))
end
\end{lstlisting}

The view sees the proxy, not the source model:

\begin{lstlisting}[style=crystal]
layers_proxy = Qt6::SortFilterProxyModel.new(main)
layers_proxy.source_model = layer_model
layers_proxy.sort_role = Qt6::ItemDataRole::User
layers_proxy.sort

layer_tree = Qt6::TreeView.new
layer_tree.model = layers_proxy
layer_tree.expand_all

layer_selection = Qt6::ItemSelectionModel.new(layers_proxy, layer_tree)
layer_tree.selection_model = layer_selection
\end{lstlisting}

When the current index changes, the example reads the selected layer name through the proxy and applies it to the state object. The canvas and inspector then refresh from state.

\section{Dock The Layer Panel}

A dock is just a titled container. Build a small widget for the dock body, then assign it to \api{DockWidget#widget}.

\begin{lstlisting}[style=crystal]
layer_summary = Qt6::Label.new("Active layer: Terrain")

layers_panel = Qt6::Widget.new
layers_panel.vbox do |column|
  column << Qt6::Label.new("Layer Manager")
  column << layer_tree
  column << layer_summary
end

layers_dock = Qt6::DockWidget.new("Layers", main)
layers_dock.widget = layers_panel
main.add_dock_widget(layers_dock, Qt6::DockArea::Left)
\end{lstlisting}

This is a good pattern for desktop tools: the central widget handles the primary work, and docks provide secondary views or controls without taking over the main surface.

\section{Build The Inspector}

The inspector is ordinary form UI: spin boxes for numeric state, a checkbox for grid visibility, buttons for commands, and a label explaining the relationship to the canvas.

\begin{lstlisting}[style=crystal]
zoom_spin = Qt6::DoubleSpinBox.new
zoom_spin.set_range(0.5, 3.0)
zoom_spin.single_step = 0.1
zoom_spin.value = state.zoom

spacing_spin = Qt6::SpinBox.new
spacing_spin.set_range(24, 96)
spacing_spin.value = state.grid_spacing

show_grid = Qt6::CheckBox.new("Show grid")
show_grid.checked = state.show_grid

accent_button = Qt6::PushButton.new("Accent Color")
reset_button = Qt6::PushButton.new("Reset View")
export_button = Qt6::PushButton.new("Export PNG")
\end{lstlisting}

Each control writes to state and schedules a repaint:

\begin{lstlisting}[style=crystal]
zoom_spin.on_value_changed do |value|
  state.zoom = value
  canvas.update
  status_bar.show_message("Zoom #{value.round(2)}x")
end

show_grid.on_toggled do |checked|
  state.show_grid = checked
  canvas.update
  status_bar.show_message(checked ? "Grid enabled" : "Grid hidden")
end
\end{lstlisting}

The right dock is assembled the same way as the layer dock.

\begin{lstlisting}[style=crystal]
inspector = Qt6::Widget.new
inspector.form do |form|
  form.add_row("Zoom", zoom_spin)
  form.add_row("Grid spacing", spacing_spin)
  form.add_row(show_grid)
  form.add_row(accent_button)
  form.add_row(reset_button)
  form.add_row(export_button)
end

inspector_dock = Qt6::DockWidget.new("Inspector", main)
inspector_dock.widget = inspector
main.add_dock_widget(inspector_dock, Qt6::DockArea::Right)
\end{lstlisting}

\begin{figure}[tbp]
  \centering
  \begin{minipage}{0.46\textwidth}
    \centering
    \textbf{Layer dock}\par\smallskip
    \screenshotimage{docs/book/images/putting-together-layers-dock.png}{width=0.96\linewidth,height=0.34\textheight,keepaspectratio}%
  \end{minipage}\hfill
  \begin{minipage}{0.48\textwidth}
    \centering
    \textbf{Inspector dock}\par\smallskip
    \screenshotimage{docs/book/images/putting-together-inspector-dock.png}{width=0.96\linewidth,height=0.34\textheight,keepaspectratio}%
  \end{minipage}
  \caption{The left dock selects scene state; the right dock edits values that repaint the canvas.}
\end{figure}

\section{Use Dialogs At The Edges}

Dialogs work best at workflow boundaries: choosing a file, choosing a color, confirming a destructive operation, or showing progress. The vertical slice uses a color dialog to change the layer accent.

\begin{lstlisting}[style=crystal]
accent_button.on_clicked do
  color = Qt6::ColorDialog.get_color(
    main,
    current_color: state.accent,
    title: "Choose Accent Color",
    show_alpha_channel: false
  )

  if color
    state.accent = color
    canvas.update
    status_bar.show_message("Accent color changed", 1400)
  end
end
\end{lstlisting}

It also uses a save-file dialog for PNG export. The export itself is just a widget grab.

\begin{lstlisting}[style=crystal]
export_png = -> do
  suggested = File.join(Dir.current, "editor-slice.png")
  output = Qt6::FileDialog.get_save_file_name(
    main,
    title: "Export Canvas Snapshot",
    directory: suggested,
    filter: "PNG Images (*.png);;All Files (*)"
  )

  if output && !output.empty?
    output = "#{output}.png" unless output.downcase.ends_with?(".png")
    canvas.grab.save(output)
  end
end
\end{lstlisting}

The complete example keeps a reusable \api{FileDialog} instance, preselects the last export path from \api{QSettings}, appends \api{.png} when needed, checks the return value from \api{canvas.grab.save}, and stores the successful path for the next export. That makes the export command a complete user workflow rather than just a screenshot helper.

The broader showcase example adds image loading through \api{QImageReader}. It also demonstrates clipboard MIME payloads, settings persistence, text editing, and drag/drop.

\section{Promote Commands To Actions}

Once a command exists as a lambda or method, promote it to an \api{Action}. Actions can live in menus and toolbars at the same time.

\begin{lstlisting}[style=crystal]
export_action = Qt6::Action.new("Export PNG", main)
export_action.shortcut = "Ctrl+E"
export_action.on_triggered { export_png.call }

reset_view_action = Qt6::Action.new("Reset View", main)
reset_view_action.shortcut = "0"
reset_view_action.on_triggered { reset_button.click }

undo_action = undo_stack.create_undo_action(main, "Undo")
redo_action = undo_stack.create_redo_action(main, "Redo")

file_menu = main.menu_bar.add_menu("File")
file_menu << export_action

edit_menu = main.menu_bar.add_menu("Edit")
edit_menu << undo_action
edit_menu << redo_action

view_menu = main.menu_bar.add_menu("View")
view_menu << reset_view_action

toolbar = Qt6::ToolBar.new("Editor", main)
toolbar << export_action
toolbar << undo_action
toolbar << redo_action
toolbar << reset_view_action
main.add_tool_bar(toolbar)
\end{lstlisting}

This is the moment where a prototype begins to feel like a desktop application. The same command can be discovered from the menu, used quickly from the toolbar, and triggered from the keyboard.

\section{Keep The Refresh Path Small}

The example uses a small refresh helper so callbacks do not duplicate status-bar and repaint behavior.

\begin{lstlisting}[style=crystal]
refresh_canvas = ->(message : String?) do
  canvas.update
  if text = message
    status_bar.show_message(text, 1800)
  end
end
\end{lstlisting}

That helper is intentionally modest. It does not hide the application architecture; it only removes repeated glue. As the app grows, this kind of helper can become a presenter object, document controller, or command layer.

\section{Where Undo Fits}

The vertical slice wires a small \api{UndoStack} around the changes that feel like editor commands: selecting a layer profile, changing inspector values, resetting the view, and choosing a new accent color. Each command captures a before and after snapshot, then lets \api{redo} and \api{undo} restore the appropriate state.

\begin{itemize}
\item Inspector controls push undoable commands instead of mutating the state directly.
\item The window title shows a dirty marker until the stack is marked clean.
\item The menu and toolbar share the stack's generated undo and redo actions.
\item Settings persistence keeps the active layer, view, grid, marker size, and export path between launches.
\item A clipboard action copies text, HTML, image, and custom MIME data for the current canvas snapshot.
\end{itemize}

The rule from Chapter 9 still applies: capture old and new state, push a command, and let \api{redo} perform the forward edit. Later features, such as moving scene objects or dropping external image assets, should follow the same command shape.

\section{Second-Pass Features}

After the first vertical slice is stable, the next natural pass is to make the editor feel more like a document tool. Keep these features small and command-shaped:

\begin{itemize}
\item Editable layer rows: make the layer model return \api{ItemFlag::Editable}, commit through \api{model_set_data}, and push a \api{Rename layer} command that updates state, the model, and the canvas together.
\item Drag/drop layer ordering: enable internal moves on the layer view, move the backing model rows with \api{begin_move_rows} and \api{end_move_rows}, and record a \api{Move layer} undo command.
\item Object movement commands: during mouse drag, update only transient preview state; when the release event arrives, push one \api{Move marker} command from the starting scene position to the final scene position.
\item External drops: accept image or text MIME payloads on the canvas, decode or store them at the document boundary, then push an import command rather than mutating the document directly from the drop callback.
\end{itemize}

\section{From Example To App}

Before turning this vertical slice into a longer-lived application, walk through the seams that prototypes usually skip:

\begin{itemize}
\item Document state: decide which fields belong to the saved document, which are view preferences, and which are transient UI state.
\item Dirty tracking: use \api{UndoStack#set_clean} only after a successful save, update the window title and Save action from \api{clean?}, and reset the stack when a different document is loaded.
\item Command enablement: disable export, copy, delete, move, or property actions when there is no valid selection or no document.
\item Save and export paths: keep project-save paths separate from PNG export paths, remember successful paths in \api{QSettings}, and check every \api{save} return value.
\item Teardown and release checks: let parented widgets be owned by Qt, release explicit resources that are not parented, and make sure painter targets are finished before reading or saving them.
\end{itemize}

\section{A Reasonable Build Order}

\begin{enumerate}
\item Build the main window, central canvas, and status bar.
\item Draw a static scene from an explicit state object.
\item Add pan and zoom callbacks that mutate state and call \api{update}.
\item Add a layer model, proxy, tree view, and selection model.
\item Put the layer view in a left dock.
\item Add inspector controls in a right dock.
\item Promote repeated commands into \api{Action} objects.
\item Add file/color dialogs at workflow boundaries.
\item Add export through \api{Widget#grab}.
\item Add undo, persistence, clipboard payloads, drag/drop, and richer document behavior after the loop above feels solid.
\end{enumerate}

\section{Reference Examples}

\api{examples/editor_vertical_slice.cr} is the most direct worked example. It is the one to read when you want to see the whole vertical slice in one file.

\api{examples/desktop_editor_showcase.cr} is broader. It demonstrates application metadata, settings, model/view panels, image loading, clipboard MIME data, text editing, drag/drop, custom preview rendering, and PNG export.

For release documentation, prefer screenshots from maintained examples over mockups. The example programs are small enough to keep fresh as the binding evolves, and concrete screenshots make later chapters easier for readers to trust.

\section{Reader Exercise}

Add one new layer property, such as route thickness or label density. Put the value in \api{EditorVerticalSliceState}, include it in \api{EditorVerticalSliceSnapshot}, add an inspector control, repaint the canvas from the new value, and make the change undoable. Then export a PNG and confirm that the captured canvas reflects the property after redo, undo, and redo again.
